\chapter{Introduction and overview}\label{ch:c1}\label{ch:intro}
This chapter explains what the coupling does, how its parts fit together, and what a Raven user needs to know before
starting.

\section{The problem this solves}
Raven \cite{craig2020raven} simulates the water cycle of a watershed: rain and snow, soil moisture, evaporation,
runoff and flow in rivers and lakes. It divides a watershed into subbasins, and each subbasin into Hydrologic
Response Units (HRUs) --- areas that behave alike (for example, forest on till, or a wetland). Raven describes
groundwater with simple ``buckets'' (linear reservoirs): water enters, and leaves at a rate proportional to what is
stored. Buckets are fast and easy to calibrate, but they cannot say \emph{where} groundwater flows, how deep the water
table is at a well, how a pumping well lowers nearby streamflow, or whether a stream gains water from the aquifer or
loses water to it.

MODFLOW~6 \cite{langevin2017mf6} is the U.S. Geological Survey's groundwater model and the international standard for
those questions. It divides the underground into a three-dimensional grid of cells and calculates the water level
(the \emph{head}) in every cell from the physics of flow through soil and rock. But MODFLOW has no snow, soil or
evaporation processes of its own; it needs to be told how much water reaches the water table.

The coupling joins the two so that each does what it is good at: Raven calculates everything at and near the land
surface; MODFLOW calculates the saturated groundwater below; and the two exchange water every day.

\begin{plain}
Raven decides how much water soaks down to the water table. MODFLOW moves that water sideways and downwards through
the aquifer. Wherever the aquifer overflows --- into rivers, wetlands, lakes or back to the soil --- MODFLOW hands the
water back to Raven. Pumping wells and other stresses act on MODFLOW's side. Both models keep exact accounts, so not
a drop is created or lost at the boundary between them.
\end{plain}

\section{Two ways to use the coupling}\label{sec:twoways}
Most users are in one of two situations (Figure~\ref{fig:twoways}).
\begin{itemize}
\item \textbf{You already have a MODFLOW~6 model} --- perhaps one calibrated over years. Raven connects to it and uses it
      as it is: same grid, same layers, same wells and boundaries (Chapter~\ref{ch:existing}).
\item \textbf{You have a Raven model but no groundwater model.} Raven builds one for you from your HRUs and a short
      description of the soil and rock layers under them (Chapter~\ref{ch:c3}).
\end{itemize}
Everything else --- how water is exchanged, how the accounts are kept, the outputs and the tests --- is the same in both cases.

\begin{figure}[!htbp]\centering% Which chapter to start from (TikZ; text is set by LaTeX, so boxes always fit their text)
\begin{tikzpicture}[font=\sffamily\small,
  base/.style={rounded corners=3pt,line width=0.8pt,align=center,inner sep=5pt,execute at begin node={\hyphenpenalty=10000\exhyphenpenalty=10000}},
  q/.style={base,draw=brand,fill=white,text width=58mm,font=\sffamily\bfseries},
  yesA/.style={base,draw=teal!70!black,fill=teal!7,text width=66mm},
  noA/.style={base,draw=orange!80!black,fill=orange!7,text width=66mm},
  yesB/.style={base,draw=teal!70!black,fill=white,text width=66mm,font=\sffamily\footnotesize},
  noB/.style={base,draw=orange!80!black,fill=white,text width=66mm,font=\sffamily\footnotesize},
  fin/.style={base,draw=brand,fill=brandlight,text width=144mm},
  arr/.style={-{Stealth[length=2.4mm]},line width=0.9pt},
  lab/.style={font=\sffamily\bfseries\small,inner sep=2pt}]
\node[q] (q) at (0,0) {Do you already have a MODFLOW~6 groundwater model?};
\node[yesA,anchor=north] (y) at (-39mm,-11mm) {\textbf{Raven connects to your model}\\(Chapter~\ref{ch:existing})\\[3pt]Your model is used as it is. Raven links its HRUs to the model's top layer.};
\node[noA,anchor=north] (n) at (39mm,-11mm) {\textbf{Raven builds a groundwater model}\\(Chapter~\ref{ch:c3})\\[3pt]From your HRUs and a short description of the soil and rock layers under them.};
\node[yesB,anchor=north] (y2) at ([yshift=-5mm]y.south) {Tell Raven where the model is on the map: its projection (\texttt{:MF6CRS EPSG:\dots}), a file of cell outlines, or a table of overlap weights.};
\node[noB,anchor=north] (n2) at ([yshift=-5mm]n.south) {Choose a grid, or keep the default: regular, rotated, quadtree, cells that follow HRU boundaries, or your own cells.};
\path let \p1=(y2.south), \p2=(n2.south) in coordinate (m) at (0,{min(\y1,\y2)-7mm});
\node[fin,anchor=north] (e) at (m) {Then, every day of the run, Raven sends recharge down and MODFLOW sends seepage and river water back up. Both keep exact water accounts (Chapters~\ref{ch:c4} and~\ref{ch:c5}).};
\draw[arr,teal!70!black] (q.west) -| node[lab,pos=0.35,above,text=teal!70!black]{Yes} (y.north);
\draw[arr,orange!80!black] (q.east) -| node[lab,pos=0.35,above,text=orange!80!black]{No} (n.north);
\draw[arr,teal!70!black] (y.south) -- (y2.north);
\draw[arr,orange!80!black] (n.south) -- (n2.north);
\draw[arr,teal!70!black] (y2.south) -- (y2.south |- e.north);
\draw[arr,orange!80!black] (n2.south) -- (n2.south |- e.north);
\end{tikzpicture}

\caption{Which chapter to start from: connecting Raven to an existing MODFLOW~6 model, or letting Raven build one.}\label{fig:twoways}
\end{figure}

\section{How to read this manual}
\begin{table}[!htbp]\centering\small
\caption{Where to find what.}\label{tab:readguide}
\begin{tabularx}{\linewidth}{P{0.34\linewidth}Y}\toprule\hdr
If you want to\dots & Read\\\midrule
understand the idea & this chapter and Chapter~\ref{ch:c2} (groundwater in plain terms)\\
connect Raven to your MODFLOW model & Chapter~\ref{ch:existing}; example \cmd{Liard\_existing\_model}\\
let Raven build a groundwater model & Chapter~\ref{ch:c3}; example \cmd{Liard\_groundwater}\\
know what water moves where, and how & Chapters~\ref{ch:c4} and~\ref{ch:c5}\\
look up a command or an output column & Chapter~\ref{ch:c6} and Appendix~\ref{app:tables}\\
set up, calibrate or fix a model & Chapter~\ref{ch:c7}\\
see how the code was tested & Chapter~\ref{ch:c8}\\
change the code & Chapter~\ref{ch:c9}\\\bottomrule
\end{tabularx}
\end{table}
Every chapter starts with a box \emph{In plain words} that says what it covers without technical terms; technical words
are explained in the glossary below.

\section{What the user does --- and what Raven does}
The user works only with Raven's familiar text files:
\begin{enumerate}
\item marks which HRUs have groundwater underneath (\cmd{AQUIFER\_PROFILE} in the \cmd{.rvh} file);
\item describes the aquifer materials and their layering (\cmd{.rvp} file);
\item gives the map of HRUs (a GeoJSON file) and, optionally, land-surface and bedrock maps, rivers, wells,
monitoring wells and boundaries (a new \cmd{.rvg} file);
\item switches the coupling on (\cmd{:GroundwaterModel MODFLOW6} in the \cmd{.rvi} file).
\end{enumerate}
Raven then does everything else automatically: it lays a grid over the HRUs, works out how much of each HRU lies in
each grid cell, builds the layers from the land surface down to bedrock, writes a standard MODFLOW~6 model, starts
MODFLOW, exchanges water every time step, and writes results next to its usual outputs. MODFLOW is used as a
\emph{library} (a program part that another program can call); Raven talks to it through MODFLOW's official
programming interface \cite{hughes2022xmi}. No MODFLOW source code is changed.

\section{Main features}
\begin{itemize}
\item Any number of \emph{aquifer profiles} (layer stacks), each with any number of layers; layer types aquifer,
confined aquifer, aquitard (confining layer) and aquiclude.
\item Land surface and bedrock from raster maps; layers that thin out against bedrock are handled automatically.
\item Coupling to an existing MODFLOW 6 model: Raven links its HRUs to the model's top active layer and leaves the model
      intact, with built-in map projections, unit conversion and safeguards against double counting (Chapter~\ref{ch:existing}).
\item Five grid options: the default regular grid, rotated and refined regular grids, quadtree grids, meshes whose cells
      follow HRU boundaries, and imported cells; layers can be connected by elevation (Section~\ref{sec:gridtypes}).
\item Water exchanges: recharge (with optional delay), groundwater seepage at the land surface, rivers (gaining and
losing), evaporation from a shallow water table, wetlands and lakes held in Raven stores, Raven reservoirs, pumping and
injection wells, head-dependent and fixed-head boundaries.
\item Cells can dry out and rewet; areas of aquifer can disconnect and reconnect by themselves.
\item Monitoring wells: simulated heads are written out and scored against observations like streamflow, so the
groundwater parameters can be calibrated (for example with Ostrich).
\item Hotstart (continuing a run from where another ended), ensembles, NetCDF output of the whole head field.
\item Exact water accounting on both sides, checked automatically by an audit tool supplied with the code.
\end{itemize}

\section{Glossary}
\begin{longtable}{P{0.24\linewidth}P{0.708\linewidth}}
\caption{Glossary.}\\
\toprule\hdr
Term & Meaning\\\midrule\endfirsthead
\caption[]{Glossary (continued).}\\
\toprule\hdr
Term & Meaning\\\midrule\endhead
Active cell & A cell that MODFLOW calculates. Cells outside the aquifer, or switched off in the model, are inactive.\\
API, library & MODFLOW~6 comes as a program library (\cmd{libmf6}) that Raven loads and controls step by step, instead of running it as a separate program.\\
Aquiclude & A layer that, for practical purposes, does not let water through.\\
Aquifer & A soil or rock layer that holds and passes on useful amounts of groundwater (gravel, sand).\\
Aquitard & A layer that lets water through only slowly (clay, till); water leaks through it.\\
Auxiliary variable & An extra column of numbers attached to each entry of a MODFLOW package --- here, for example, the Raven subbasin that a river cell belongs to.\\
Bedrock & The base of the groundwater system in this model; below it no flow is simulated.\\
Cell & One box of the MODFLOW grid. A \emph{column} is the stack of cells at one map location.\\
Conductance & How easily water passes a connection, in m$^2$/d; flow = conductance $\times$ difference in water level.\\
Confined / unconfined & An unconfined aquifer has a free water table; a confined one is completely full and under pressure.\\
Converged & MODFLOW solves each day by repeated improvement; a day has \emph{converged} when the answer stops changing within a small tolerance.\\
Coupled HRU & An HRU that has groundwater underneath it in the model and so exchanges water with MODFLOW.\\
DIS, DISV, DISU & MODFLOW's three ways of describing a grid: rows and columns; layers of any-shaped cells; fully free cells.\\
EPSG code & A standard number for a map projection, e.g.\ EPSG:32610 is UTM zone 10 North. Used with \cmd{:MF6CRS}.\\
Existing-model mode & Using the coupling with a MODFLOW~6 model you already have (Chapter~\ref{ch:existing}).\\
Grid & The set of cells MODFLOW calculates in: squares of one size by default, or other shapes (Section~\ref{sec:gridtypes}).\\
Head ($h$) & The level to which water rises in a well, in metres above sea level. Water flows from high to low head.\\
Hotstart, restart & Starting a run from the saved end state (\cmd{solution.rvc}) of an earlier run, so a long run can be split or continued.\\
HRU & Hydrologic Response Unit: Raven's basic land unit, an area that behaves alike.\\
Hydraulic conductivity ($K$) & How easily water moves through a material, in m/d.\\
Leakance & How easily water passes a thin layer (a riverbed, the ground surface), per square metre, in 1/d.\\
Linkage, overlap & How much of each HRU lies in each grid cell. It decides where each HRU's water goes.\\
Newton method & The way MODFLOW solves its equations here; it copes well with cells drying out and filling up again.\\
Overlap weight & The share of a cell covered by an HRU (0 to 1); a table of them can replace the map of cells (\cmd{:OverlapWeights}).\\
Package & One component of a MODFLOW model, in its own file: wells (WEL), rivers (RIV), recharge (RCH) and so on.\\
Pass-through cell & A very thin cell where a layer is missing; water passes straight through it from the layer above to the layer below.\\
Profile & A stack of aquifer layers assigned to HRUs (used when Raven builds the model).\\
Projection & The rule that turns positions on the round Earth into flat map coordinates in metres; MODFLOW models use one, usually without saying which.\\
Quadtree, mesh & Grids with cells of different sizes: a quadtree cuts squares into four where detail is wanted; a mesh (here) has cells shaped to follow HRU boundaries.\\
Recharge & Water that reaches the water table from above.\\
Seepage & Groundwater flowing out at the land surface where the water table reaches it.\\
Specific storage ($S_s$) & Water released per m$^3$ of aquifer per metre of head fall when the aquifer stays full [1/m].\\
Specific yield ($S_y$) & Water released per m$^2$ of water table per metre of fall when the aquifer drains [--].\\
Stage & Water level in a river, lake or reservoir [m].\\
Steady state & A calculation as if conditions never changed --- the long-term balance. Often used for the first period of a model.\\
Stress period & A stretch of time in a MODFLOW model (often a month) during which inputs such as pumping rates stay the same.\\
Time step & One step of the calculation; in the coupling, one Raven time step (usually a day).\\
Transmissivity ($T$) & $K$ times saturated thickness, in m$^2$/d.\\
Water table & The top of the saturated zone in an unconfined aquifer.\\
Working copy & The copy of an existing MODFLOW model that Raven runs, in \cmd{output/mf6/}; your own files are never changed.\\
XT3D & An optional, more exact way for MODFLOW to calculate flow between cells of odd shapes; slower and off by default.\\\bottomrule
\end{longtable}

\section{The big picture}\label{ch:overview}
\subsection{Who owns which water}
Every drop of water in the coupled system belongs to exactly one of the two models at any time:
\begin{itemize}
\item \textbf{Raven} owns snow, canopy, soil and vadose-zone stores, depressions and wetlands, rivers, lakes and
reservoirs.
\item \textbf{MODFLOW} owns the saturated groundwater below the \emph{model top}. By default the model top is the land
surface (as in GSFLOW \cite{markstrom2008gsflow}); users with physically based soil layers can lower it with
\cmd{SOIL\_ZONE\_DEPTH}.
\end{itemize}
Water moves between the two only through the exchanges described in Chapter~\ref{ch:c4}, and each exchange is recorded on both
sides. In a coupled HRU, Raven's \cmd{GROUNDWATER} store is a \emph{pass-through}: whatever enters it during a day is
handed to MODFLOW as recharge that day.

\subsection{The pieces of the software}
Figure~\ref{fig:arch} shows the parts. The \emph{model builder} runs once at the start: it reads the inputs and writes a
standard MODFLOW~6 model. The \emph{exchange} runs every time step.

\begin{figure}[!htbp]\centering
\begin{tikzpicture}[box/.style={draw,rounded corners,align=center,minimum height=11mm,text width=34mm,font=\small},
 ar/.style={-{Stealth[length=2.2mm]},thick}]
\node[box,fill=teal!12] (in) at (0,1.1) {Raven inputs\\ \cmd{rvi rvh rvp rvg rvc rvt}};
\node[box,fill=gray!12] (geo) at (0,-1.1) {Maps: HRU polygons, DEM, bedrock, rivers};
\node[box,fill=gray!20] (bld) at (5,0) {Model builder\\ grid, layers, links,\\ MODFLOW 6 files};
\node[box,fill=violet!12] (mf) at (10,0) {\cmd{libmf6}\\ MODFLOW 6};
\node[box,fill=teal!12] (rv) at (10,-3.2) {Raven solver\\ HRU processes, routing};
\node[box,fill=teal!6] (out) at (5,-3.2) {Outputs\\ diagnostics, hotstart,\\ budget, heads};
\draw[ar] (in) -- (bld); \draw[ar] (geo) -- (bld);
\draw[ar] (bld) -- node[above,font=\scriptsize]{once} (mf);
\draw[ar,<->] (mf) -- node[right,font=\scriptsize,align=left]{water and heads\\ every time step} (rv);
\draw[ar] (rv) -- (out);
\end{tikzpicture}
\caption{The software pieces. The builder runs once per simulation (or per ensemble member); the exchange runs every
time step.}\label{fig:arch}
\end{figure}

\subsection{One day in the life of the coupled model}
The following happens every time step (one day in the examples); Section~\ref{ch:timestep} gives the exact algorithm.
\begin{enumerate}
\item \textbf{Raven's land-surface day.} Raven calculates snowmelt, infiltration, evaporation and flow between its soil
layers. In coupled HRUs, water that percolates past the soil ends up in \cmd{GROUNDWATER}.
\item \textbf{Recharge handed over.} The water in each HRU's \cmd{GROUNDWATER} is spread over the grid cells under that
HRU, in proportion to how much of the HRU lies in each cell.
\item \textbf{Stresses set.} River levels (from Raven's rivers at the start of the day), well pumping rates,
boundary heads, wetland and lake levels, and the evaporation demand not met by Raven are passed to MODFLOW.
\item \textbf{MODFLOW's day.} MODFLOW calculates new heads everywhere, including the flows to and from rivers,
drains, wells and boundaries.
\item \textbf{Water handed back.} Groundwater that seeps out at the land surface is added to a Raven store of the HRUs
above; groundwater that discharges to rivers is added to the river inflow of the subbasin; river losses are taken from
the river; wetland and lake exchanges change those Raven stores.
\item \textbf{Raven's river day.} Raven routes water down the river network, now including the groundwater
contributions. Reservoirs compute their seepage using the aquifer head beneath them.
\item \textbf{Accounts.} Both models update their water balances, and the groundwater outputs are written.
\end{enumerate}

\subsection{Design principles}
\begin{enumerate}
\item \textbf{Raven is the single user interface.} All inputs are Raven files; MODFLOW files are generated.
\item \textbf{Standard MODFLOW only.} Only standard MODFLOW~6 packages are used; the generated model can be opened in
FloPy or ModelMuse.
\item \textbf{Each fact is entered once.} Geometry comes from maps, hydrogeology from profiles, time-varying stresses
from time series.
\item \textbf{Generic.} Nothing is specific to a basin; any Raven model can couple any subset of its HRUs.
\item \textbf{Exact accounting.} Every exchanged volume is computed once and applied identically on both sides.
\item \textbf{No surprise for existing users.} Models without groundwater give bit-for-bit the same results as
Raven~4.15 (the version this coupling is built on); Raven still builds and runs without MODFLOW installed.
\end{enumerate}

\section{Scope and assumptions}
\begin{itemize}
\item \textbf{What is simulated by MODFLOW:} saturated groundwater flow (Darcy flow, storage change) in a layered aquifer
system under the coupled HRUs, including dry and rewetting cells.
\item \textbf{What stays in Raven:} everything above the model top --- snow, canopy, soil and unsaturated zone, surface
runoff, depressions, rivers, lakes and reservoirs --- and all HRUs without an aquifer profile.
\item \textbf{Water density and temperature} are constant; there is no salt-water or heat transport.
\item \textbf{The grid} is regular (square cells); the model area is the part of the rectangle around the coupled HRUs
that those HRUs cover.
\item \textbf{Time}: both models use Raven's time step (daily or shorter). Exchanges are explicit between the models at
the step level, but every head-dependent exchange is implicit inside MODFLOW.
\end{itemize}

\section{When to use the coupling --- and when not to}
Use it when the question involves the \emph{spatial} behaviour of groundwater: water-table depth at a location, effects of
pumping on streams, gaining and losing reaches, groundwater flow between subbasins, wetlands fed by groundwater, drought
recession of large aquifers, or when head measurements should constrain the model. Raven's conceptual baseflow remains the
better choice for fast screening, very large ensembles, or basins without a significant aquifer.

\section{What changes for an existing Raven model}
\begin{itemize}
\item Nothing, unless \cmd{:GroundwaterModel MODFLOW6} is added: results stay bit-for-bit identical.
\item With the coupling on, coupled HRUs lose their conceptual deep-groundwater behaviour (their \cmd{GROUNDWATER} store
becomes a pass-through), and their baseflow now comes from the aquifer via seepage and river exchange.
\item Run times increase: the 20-year Liard run (3\,252 active columns) takes about 15 minutes instead of 35 seconds on one processor core.
\end{itemize}

\section{Files at a glance}
\begin{table}[!htbp]\centering
\caption{Input and output files of the coupling.}
\begin{tabular}{P{0.2\linewidth}P{0.748\linewidth}}\toprule\hdr
File & Role in the coupling\\\midrule
\cmd{.rvi} & switch on (\cmd{:GroundwaterModel MODFLOW6}); route recharge into \cmd{GROUNDWATER}\\
\cmd{.rvh} & \cmd{AQUIFER\_PROFILE} per HRU; reservoirs with \cmd{:SeepageParameters}\\
\cmd{.rvp} & aquifer classes, profiles and profile parameters\\
\cmd{.rvg} & geometry (HRU polygons, maps, rivers), grid, wells, monitoring wells, boundaries, options\\
\cmd{.rvt} & well rates, boundary heads, observed heads\\
\cmd{.rvc} & initial heads; hotstart blocks written by Raven\\
GeoJSON, ASCII grids & HRU polygons, river lines, boundary shapes; DEM and bedrock maps\\
\cmd{<output>/mf6/} & generated MODFLOW 6 model and its own outputs\\\bottomrule
\end{tabular}
\end{table}
